\subsection{Before and after re-ranking}
\label{sec:beforeafter}

\begin{table*}[t]
\centering\scriptsize
\begin{tabular}{@{}lccccc@{}}
\toprule
system & AUC & MRR & MRR, all clicks & nDCG@5 & nDCG@10 \\
\midrule
\multicolumn{6}{@{}l}{\textit{\eb{} small, test}} \\
BM25 (stage one) & \cellci{0.5107}{0.5094}{0.5120} & \cellci{0.3257}{0.3245}{0.3269} & \cellci{0.3253}{0.3241}{0.3264} & \cellci{0.3577}{0.3563}{0.3591} & \cellci{0.4409}{0.4397}{0.4420} \\[2pt]
embeddings (stage one) & \cellci{0.5397}{0.5384}{0.5409} & \cellci{0.3492}{0.3480}{0.3503} & \cellci{0.3487}{0.3475}{0.3498} & \cellci{0.3831}{0.3818}{0.3845} & \cellci{0.4627}{0.4616}{0.4638} \\[2pt]
fused (stage one) & \cellci{0.5380}{0.5367}{0.5392} & \cellci{0.3474}{0.3463}{0.3486} & \cellci{0.3470}{0.3458}{0.3481} & \cellci{0.3813}{0.3800}{0.3827} & \cellci{0.4613}{0.4602}{0.4624} \\[2pt]
two-stage re-ranker & \cellcib{0.7908}{0.7898}{0.7918} & \cellcib{0.5685}{0.5671}{0.5699} & \cellcib{0.5678}{0.5665}{0.5693} & \cellcib{0.6330}{0.6318}{0.6344} & \cellcib{0.6650}{0.6639}{0.6662} \\[2pt]
fused + lifetime popularity$^\ast$ & \cellci{0.5797}{0.5785}{0.5810} & \cellci{0.3688}{0.3677}{0.3700} & \cellci{0.3684}{0.3672}{0.3696} & \cellci{0.4101}{0.4088}{0.4115} & \cellci{0.4841}{0.4830}{0.4852} \\[2pt]
\midrule
\multicolumn{6}{@{}l}{\textit{\mind{} small, test}} \\
BM25 (stage one) & \cellci{0.5685}{0.5663}{0.5706} & \cellci{0.3108}{0.3085}{0.3133} & \cellci{0.2693}{0.2673}{0.2717} & \cellci{0.2868}{0.2844}{0.2894} & \cellci{0.3479}{0.3456}{0.3504} \\[2pt]
embeddings (stage one) & \cellci{0.6369}{0.6347}{0.6390} & \cellci{0.3510}{0.3486}{0.3535} & \cellci{0.3061}{0.3039}{0.3085} & \cellci{0.3341}{0.3316}{0.3367} & \cellci{0.3937}{0.3913}{0.3961} \\[2pt]
fused (stage one) & \cellci{0.6381}{0.6358}{0.6402} & \cellci{0.3536}{0.3511}{0.3561} & \cellci{0.3075}{0.3053}{0.3098} & \cellci{0.3361}{0.3335}{0.3388} & \cellci{0.3956}{0.3934}{0.3981} \\[2pt]
two-stage re-ranker & \cellcib{0.6964}{0.6943}{0.6983} & \cellcib{0.4027}{0.4001}{0.4054} & \cellcib{0.3459}{0.3436}{0.3483} & \cellcib{0.3854}{0.3827}{0.3882} & \cellcib{0.4458}{0.4435}{0.4483} \\[2pt]
\bottomrule
\end{tabular}

\caption{All accuracy metrics on the test split, with 95\% bootstrap intervals beneath each mean.
``Stage one'' rows are the retrieval scores ranking the same candidates; the re-ranker reorders
them. $^\ast$Lifetime popularity is not available at serving time and is reported only to price
what it would buy (Section~\ref{sec:q9}). Validation-split figures are in Appendix~\ref{app:val}.}
\label{tab:overall}
\end{table*}

\textbf{Two baselines appear in this report and they are not the same thing.} Here ``stage one''
is the fused retrieval score ranking the candidates directly (\EBFusedTest{} on test). The
ablation grid in Table~\ref{tab:choices} instead compares against \feat{stage1_only}, a ranker
trained on the two retrieval scores alone, which is slightly stronger at 0.5430. That is why the
same result reads \EBGainTest{} here and $+$0.2223 there.

\textbf{On \eb{} the re-ranker is worth \EBGainTest{} AUC} over fused retrieval,
\EBGainTestCI{} on test, and \EBGainVal{} \EBGainValCI{} on val. The val and test intervals
overlap, so the gain generalises across the time boundary rather than being an artefact of what we
selected on. Every other accuracy metric moves the same way, so this is a genuine reordering rather
than a trade between metrics.

\textbf{On \mind{} the same architecture is worth \MDGainTest{} \MDGainTestCI{}}, far less. Two
things compound. Five of the 22 features have no source in \mind{}, including every freshness and
engagement signal, so the behavioural axis is thinner. And \mind{}'s stage one is already stronger
in absolute terms (\MDFusedTest{} against \eb{}'s \EBFusedTest{}), which leaves less headroom to
begin with. We report the contrast rather than presenting one system that is quietly two: a
smaller gain on a dataset with less signal is a measured property of the data.

\textbf{One family of features did transfer, and it was the one built under duress.} An earlier
round of seven features moved \mind{} by nothing significant at all. The five exposure features,
added because Codabench withholds the click labels, are worth \LadderExpMindTest{}
\LadderExpMindTestCI{} on \mind{} test and \LadderExpMindVal{} \LadderExpMindValCI{} on val. How
often an article was \emph{shown} is apparently a signal both corpora carry, where freshness and
engagement are specific to \eb{}.
